\documentclass[11pt,a4paper]{article}
\usepackage[T1]{fontenc}
\usepackage{lmodern}
\usepackage[a4paper,margin=2.2cm]{geometry}
\usepackage{booktabs}
\usepackage{tabularx}
\usepackage{array}
\usepackage{amssymb}
\usepackage{enumitem}
\usepackage{xcolor}
\usepackage{titlesec}
\usepackage{fancyhdr}
\usepackage{xurl}
\usepackage[hidelinks]{hyperref}

\definecolor{accent}{HTML}{1F4E79}
\definecolor{done}{HTML}{2E7D32}
\definecolor{now}{HTML}{B26A00}
\titleformat{\section}{\large\bfseries\color{accent}}{\thesection.}{0.5em}{}
\titlespacing*{\section}{0pt}{1.1em}{0.4em}
\setlist{itemsep=2pt,topsep=3pt}
\renewcommand{\arraystretch}{1.25}
\setlength{\parindent}{0pt}
\setlength{\parskip}{0.45em}

\pagestyle{fancy}
\fancyhf{}
\lhead{\small Major Project: Status Report}
\rhead{\small MCP-Lock}
\cfoot{\small \thepage}
\renewcommand{\headrulewidth}{0.4pt}

\newcommand{\done}{\textcolor{done}{\checkmark~Completed}}
\newcommand{\inprog}{\textcolor{now}{$\blacktriangleright$~In progress}}

\begin{document}

\begin{center}
  {\LARGE\bfseries\color{accent} Major Project: Status Report}\\[0.5em]
  {\large SRM Institute of Science and Technology}\\[0.3em]
  {\normalsize 24 October 2026}
\end{center}
\vspace{0.4em}
\hrule
\vspace{0.6em}

\section{Team Members}
\begin{tabularx}{\textwidth}{@{}cX@{}}
\toprule
\textbf{No.} & \textbf{Name} \\
\midrule
1 & Aakash G S \\
2 & Hariharan \\
3 & Jaswanth \\
\bottomrule
\end{tabularx}

\section{Proposed Project Title}
\textbf{MCP-Lock: Lockfiles and Pre-Launch Gating for Coding-Agent MCP Servers}

\section{Problem Statement}
AI coding agents (for example Claude Code and Cursor) get their tools from \emph{Model Context Protocol (MCP)
servers}, small programs that agents usually download and start from public package registries with commands such
as \texttt{npx -y package} or \texttt{uvx package}. These commands are \emph{unpinned}: every launch silently fetches
the newest version. Two things can therefore change between one session and the next without the user noticing:
\begin{itemize}
  \item the \textbf{code} that runs on the developer's machine, which can carry malware if a maintainer account is
        hijacked or a malicious update is published; and
  \item the \textbf{tool descriptions} that the AI model reads as instructions, which an attacker can rewrite after
        the user has approved the server (a ``rug pull'' that injects hidden instructions into the model).
\end{itemize}
Existing pre-install checks look only at package names, sources and age, and registry provenance records who built
a package, not what tools it exposes. Our census found that \textbf{36.5\%} of 2,000 sampled MCP packages changed
their tools within 90 days, and that \textbf{97.0\%} of those changes carried no warning sign in registry metadata.

\textbf{Our project solves this} with MCP-Lock: a lockfile that records the exact reviewed state of every MCP server
(version, archive digest, build provenance, publisher, install scripts and a hash of its tool definitions), and a
gate that checks every launch against it \emph{before} any package code runs, blocking or flagging anything that
changed.

\section{Current Status}
\begin{tabularx}{\textwidth}{@{}lX@{}}
\toprule
\textbf{Stage} & \textbf{Status} \\
\midrule
Discussion stage & \done \\
Topic finalization & \done \\
Literature survey & \done{} (30 cited works, citations verified) \\
Design & \done{} (lockfile specification and gate design) \\
Implementation & \done{} (working tool, 103 automated tests passing) \\
Evaluation & \done{} (pre-registered evaluation and adaptive-attacker study) \\
Research paper & \inprog{} (6-page IEEE-format draft, revised after an internal review) \\
\bottomrule
\end{tabularx}

\textbf{Overall stage: implementation and evaluation are complete; the team is now finalizing the research paper.}

\section{Work Completed So Far}
\begin{tabularx}{\textwidth}{@{}>{\raggedright\arraybackslash}p{3.1cm}X@{}}
\toprule
\textbf{Phase} & \textbf{Outcome} \\
\midrule
Literature survey and baseline analysis &
Surveyed MCP security, agent supply-chain and package supply-chain work. Re-implemented the closest existing
defence (the Bagmar and Saraf pre-install hook, arXiv:2607.15143) as our baseline. \\
Census of MCP packages &
Collected metadata for all 15,272 eligible MCP packages on npm and PyPI (228,044 releases), and statically
extracted tool definitions from a stratified sample of 2,000. Validated the extractor against live servers in an
isolated container (precision 92.2\%, recall 78.0\%). \\
Design and implementation &
Built the \texttt{mcp-lock.json} lockfile format with a written specification and JSON schema, the
\texttt{mcplock} command-line tool, and a Claude Code hook that checks every install or launch command. \\
Pre-registered evaluation &
Registered hypotheses and analysis plan on OSF before running (DOI 10.17605/OSF.IO/2AMJ3). On 490 real attacks
MCP-Lock detected \textbf{39.8\%} against \textbf{19.6\%} for the baseline (99 more attacks; exact McNemar
$p = 3.2\times10^{-30}$). It caught all 40 of 40 synthetic rug pulls, which the baseline missed entirely. \\
Adaptive-attacker study &
Built 12 attacks by an attacker who knows every check. The basic gate blocks 1 and flags 4 of them; five
optional hardening switches close all but 2, and the cost of each was measured on 8,180 honest upgrades. Ran every
test attack server in a sealed container to confirm the runtime results (11 of 11 agree). \\
Safer-alternative suggestions &
When a command is blocked, the gate now proposes a verified safe replacement (for example the last clean
version), shown only if it passes the gate. Correct for 114 of 114 hijacked packages that received a suggestion. \\
Research paper &
Wrote a 6-page IEEE-format paper in which every number is generated from the result files. Ran a simulated IEEE-style
peer review and fixed 9 of its 10 points: added comparisons with simpler alternatives, prior work on tool pinning,
and robustness checks. \\
\bottomrule
\end{tabularx}

Code, data and results are public at \url{https://github.com/ftaakash/MCP-Lock}; the pre-registration is at
\url{https://osf.io/2amj3}.

\section{Next Tasks Planned}
\begin{enumerate}
  \item \textbf{Guide review of the paper draft:} incorporate feedback from our project guide.
  \item \textbf{Select the target venue:} choose an IEEE conference or workshop, then adjust the page limit,
        anonymisation and formatting to its rules.
  \item \textbf{Expand the methods section:} describe the six documented deviations from the pre-registration in
        full, the last open point of the internal review.
  \item \textbf{Final checks before submission:} repeat the related-work search, since the field changes
        monthly; proofread; and verify every reference.
  \item \textbf{Submit the paper,} and optionally post a preprint on arXiv.
\end{enumerate}

\section{Difficulties and Support Required}
\begin{itemize}
  \item \textbf{Choice of venue and timeline:} we would value the guide's advice on a suitable IEEE conference or
        workshop and its submission deadline.
  \item \textbf{Independent evaluation:} the adaptive attacks were designed by our own team; feedback from the guide
        or another faculty member on further attacks would strengthen the evaluation.
  \item \textbf{Limited real-world data:} no public dataset of real tool-description rug pulls exists, so that part
        of the evaluation relies on synthetic and adaptive cases. We have stated this openly in the paper.
  \item \textbf{Prior work on tool pinning:} recent open-source tools already pin MCP tool definitions, so we have
        repositioned the contribution around pre-launch checking and the measured evaluation. We would welcome the
        guide's view on this framing.
\end{itemize}

\end{document}
